% ============================================================
% Theory and Methods — SCX: State-Conditioned eXpertise
% Consistency for Training Data Quality
%
% Nature Computational Science — flagship paper
% Style: lightweight, intuitive, proofs in SI
% ============================================================

\section{State-Conditioned eXpertise}
\label{sec:method}

The core challenge in automated training data cleaning is distinguishing between two fundamentally different kinds of failures: samples whose labels are wrong (noise) and samples that are intrinsically difficult but correctly labeled (hardness). Existing approaches typically conflate these two cases, treating any sample on which a model underperforms as suspicious. We show that this conflation is not a practical shortcoming but a theoretical necessity that can only be resolved with additional structure. Our method, SCX (State-Conditioned eXpertise), provides this structure through two components: a two-layer state discovery that partitions the data into semantically meaningful regimes, and a multi-expert consistency test that detects noise with provable guarantees.

\subsection{Two-Layer State Discovery}
\label{sec:two_layer}

A fundamental prerequisite for distinguishing noise from hardness is the ability to identify when two samples belong to the same \emph{state}---that is, the same region of the data manifold where model behavior is homogeneous. In the AlN machine-learned interatomic potential (MLIP) application, for example, a thermal snapshot at 900~K and an equilibrium crystal structure represent fundamentally different states: the model's expected error rate is orders of magnitude higher on the former, yet both are correctly labeled. A noise detector that does not account for this state difference would mistakenly flag all thermal snapshots as mislabeled.

We address this with a two-layer state discovery procedure. The first layer encodes domain-knowledge features derived from the raw input: for MLIP, these are Atomic Cluster Expansion (ACE) descriptors capturing the local chemical environment of each atom; for vision tasks, they are intermediate CNN latent representations; for tabular data, they are the raw feature vectors augmented with pairwise interaction terms. These features embed the domain's natural notion of similarity---atoms in similar bonding environments, or images depicting similar objects---into a vector space where distance correlates with physical or semantic proximity.

The second layer is an error-driven encoder that learns which dimensions of the feature space are most predictive of model error. The key insight is that not all feature dimensions carry equal signal about data quality. In the AlN dataset, for instance, the ACE descriptors capture both equilibrium bond geometry (which correlates weakly with model error) and thermal displacement magnitude (which correlates strongly). The error-driven encoder applies a supervised transformation that orients the feature axes to maximize the mutual information between the resulting state assignments and the observed expert disagreement rates. This is implemented as a linear projection followed by k-means clustering, with the projection matrix optimized via a surrogate objective that penalizes within-cluster variance of the expert error rates.

The effect of the second layer is substantial. Figure~\ref{fig:two_layer_improvement} shows the state refinement for AlN v3 on a subset of frames originally assigned to a single coarse state---what we call a 50\% megastate---spanning both thermal and MD configurations. The first-layer features alone (ACE descriptors clustered directly) produce a single state encompassing roughly half the dataset, achieving a noise detection F1 score of only 0.253 on this subset. After the error-driven encoder refines this megastate into six error-distinct substates, the F1 score on the same subset rises to 0.585---a 2.3-fold improvement.

\begin{figure}[t]
\centering
\includegraphics[width=\columnwidth]{figures/two_layer_improvement.pdf}
\caption{Two-layer state discovery improves noise detection F1 from 0.253 to 0.585 on a 50\% megastate in the AlN v3 dataset. The first layer (ACE descriptors) groups frames by chemical environment but conflates thermal and MD configurations that have very different error profiles. The second layer (error-driven encoder) separates these into six substates, each with homogeneous expert error rates.}
\label{fig:two_layer_improvement}
\end{figure}

The complete state discovery, applied to all 534 bulk-core AlN frames, partitions the data into eight distinct states. Notably, the noise is not uniformly distributed: all 53 detected noise frames are concentrated in the thermal (49\%) and MLMD (35\%) batches, with zero noise in equilibrium, EOS, elastic, and phonon states---consistent with the physical expectation that far-equilibrium configurations are harder to label reliably.

\subsection{Multi-Expert Consistency for Noise Detection}
\label{sec:consistency}

Once the data is partitioned into states within which expert error rates are approximately homogeneous, we can detect mislabeled samples through a simple but powerful principle: \emph{if all independently-trained experts fail on a sample, the label is likely wrong; if only some fail, the sample is likely hard but correctly labeled.}

Formally, we train $M$ expert models on disjoint subsets of the training data. For each sample $x$, each expert predicts a label $\hat{y}_m$, which we compare to the given label $y$. The consistency score $C(x)$ is the fraction of experts that fail (predict incorrectly). Our noise detection rule is: flag sample $x$ as noisy if $C(x) > \tau(s)$, where $\tau(s)$ is a state-dependent threshold calibrated to the expected expert error rate in state $s$.

The intuition is statistical. Consider a sample $x$ whose true label is $y^*$ and whose given label is $y$. There are two possible explanations for expert failure:

\begin{enumerate}
    \item \textbf{Noise}: $y \neq y^*$ (the given label is incorrect). In this case, all experts are predicted against a wrong standard. If they are independently trained and the noise is uniformly distributed, the probability that a randomly chosen expert fails is at least $1 - \rho$, where $\rho$ is the expert's clean accuracy. By independence, the probability that all $M$ experts fail simultaneously is at least $(1 - \rho)^M$, which is high for modest $M$.
    \item \textbf{Hardness}: $y = y^*$ (the label is correct) but the sample is intrinsically difficult. In this case, each expert fails independently with some probability $p < 1$. The probability that all $M$ fail is $p^M$, which vanishes exponentially as $M$ grows.
\end{enumerate}

The gap between these two regimes is what makes multi-expert consistency effective, and we can quantify it exactly.

\subsubsection{Theoretical Guarantee}

We prove the following guarantee (full proof in Supplementary Information, Section~S1). Under assumptions A1--A6 (expert training on disjoint data, conditional independence given the state, bounded loss, uniform independent noise, state-homogeneous error rates, and balanced error distribution), the noise detection F1 score satisfies:

\begin{equation}
\mathrm{F1} \;\geq\; 1 \;-\; \frac{1}{\eta} \sum_{s \in \mathcal{S}} \rho_s \exp\Bigl(-2M\,\Delta_s^2\Bigr),
\label{eq:theorem1}
\end{equation}

where $\eta$ is the global noise rate, $\rho_s$ is the proportion of data in state $s$, $M$ is the number of experts, and $\Delta_s = \min(\theta - \mu_s,\; 1 - C_{\text{bal}}\cdot\mu_s/(K-1) - \theta)$ is the \emph{separation gap} for state $s$---the minimum distance from the threshold $\theta$ to the expected consistency score under clean data ($\mu_s$) and under noise ($1 - C_{\text{bal}}\cdot\mu_s/(K-1)$). When all states share the same gap $\Delta$ (homogeneous separation), the bound simplifies to $\mathrm{F1} \geq 1 - (1/\eta)\exp(-2M\Delta^2)$.

The exponential dependence on $M$ is the key practical takeaway. For a $K$-class problem with per-state clean error rate $\mu_s$, the optimal threshold yields $\Delta_s^* = \frac{1}{2}(1 - C_{\text{bal}}\cdot\mu_s\cdot K/(K-1))$. With $K = 10$ classes (typical for image classification) and clean error rate $\mu_s = 0.15$, we obtain $\Delta^* \approx 0.42$, so that $M = 20$ experts and noise rate $\eta = 0.10$ guarantee $\mathrm{F1} \geq 0.97$. With binary classification ($K = 2$) and the same $\mu_s = 0.15$, we obtain $\Delta^* = 0.35$, requiring $M = 27$ experts to reach the same guarantee. Table~\ref{tab:f1_guarantee} shows the minimum $M$ required across practical parameter regimes; note that the Hoeffding bound used here is conservative---the Chernoff (KL-divergence) form given in the SI is typically 2--5$\times$ tighter.

\begin{table}[t]
\caption{Minimum number of experts $M$ required to guarantee $\mathrm{F1} \geq 1 - \varepsilon$ under the bound in Equation~\eqref{eq:theorem1}, for different separation gaps $\Delta$ and noise rates $\eta$. Values are computed as $M = \lceil \ln(1/(\eta\varepsilon)) \,/\, (2\Delta^2) \rceil$. The Chernoff (KL-divergence) form in the SI yields 2--5$\times$ smaller requirements; the Hoeffding values shown here are conservative upper bounds.}
\label{tab:f1_guarantee}
\begin{ruledtabular}
\begin{tabular}{cccccc}
\toprule
\multirow{2}{*}{$\varepsilon$} & \multirow{2}{*}{Target F1} & \multicolumn{2}{c}{$\eta = 0.05$} & \multicolumn{2}{c}{$\eta = 0.10$} \\
\cmidrule{3-6}
& & $\Delta = 0.25$ & $\Delta = 0.35$ & $\Delta = 0.25$ & $\Delta = 0.35$ \\
\midrule
$0.10$ & $\geq 0.90$ & 43 & 22 & 37 & 19 \\
$0.05$ & $\geq 0.95$ & 48 & 25 & 43 & 22 \\
$0.01$ & $\geq 0.99$ & 61 & 31 & 56 & 29 \\
$0.001$ & $\geq 0.999$ & 80 & 41 & 74 & 38 \\
\bottomrule
\end{tabular}
\end{ruledtabular}
\end{table}

The separation gap $\Delta$ plays a critical role. It is maximized when experts are diverse---trained on different data subsets with different random initializations---so that their errors decorrelate, and when the state discovery is accurate enough that clean-sample error rates are well-characterized. The number of classes $K$ also matters: larger $K$ increases the noise-side gap (a noisy label is less likely to coincide with an expert's mistaken prediction), yielding larger $\Delta$. This connects directly to the two-layer state discovery above: better state discovery increases $\Delta$ by making within-state error rates more homogeneous, which tightens the theoretical guarantee.

In the AlN v3 application, training $M = 12$ experts with disjoint data subsets across $K = 8$ discovered states yields per-state separation gaps ranging from $\Delta = 0.35$ (thermal state) to $\Delta = 0.47$ (equilibrium state). The global noise rate is $\eta = 53/534 \approx 0.10$. Applying the full state-weighted bound (Equation~\ref{eq:theorem1}) with the observed per-state gaps gives a theoretical F1 lower bound of approximately $0.86$ (assuming roughly equal state sizes) to $0.77$ (under conservative assumptions where the noise-heavy thermal and MLMD states are larger). The empirical $\mathrm{F1} = 0.87$ (Table~\ref{tab:aln_f1}) is consistent with the theoretical prediction, lying within or above the bound's plausible range. Per-state detection results are shown below.

\begin{table}[t]
\caption{SCX noise detection performance on the AlN v3 dataset ($M = 12$, $\eta \approx 0.10$, 8 states, 534 frames). The global theoretical F1 lower bound from Equation~\eqref{eq:theorem1} is approximately $0.77$--$0.86$ depending on state proportions; the empirical $\mathrm{F1} = 0.87$ is consistent with this range. Per-state columns show the clean error rate $\mu_s$, the separation gap $\Delta_s$, and the empirical F1.}
\label{tab:aln_f1}
\begin{ruledtabular}
\begin{tabular}{lcccc}
\toprule
State & $\mu_s$ & $\Delta_s$ & Noise frames & Empirical F1 \\
\midrule
Equilibrium & $0.03$ & $0.47$ & 0 & $1.00$ \\
EOS/strain & $0.04$ & $0.46$ & 0 & $0.97$ \\
Phonon & $0.06$ & $0.44$ & 0 & $1.00$ \\
Thermal & $0.18$ & $0.35$ & 26 (49\%) & $0.89$ \\
MLMD & $0.15$ & $0.41$ & 19 (35\%) & $0.91$ \\
3 other states & $0.04$ & $0.46$ & 8 (16\%) & $0.96$ \\
\midrule
Global & -- & -- & 53 & $\mathbf{0.87}$ \\
\bottomrule
\end{tabular}
\end{ruledtabular}
\end{table}

\subsection{When the Method Fails: Weak Feature Detection}
\label{sec:weak_features}

The guarantee in Equation~\eqref{eq:theorem1} assumes that the two-layer state discovery produces states that are genuinely informative about model error. But what if the available features are simply too weak to reveal the underlying state structure? In materials science, this arises when descriptors are too low-resolution to distinguish different bonding environments; in medical imaging, it occurs with low-resolution thumbnail images where lesion subtypes are visually indistinguishable.

We formalize this limitation. Let $\phi(x)$ be the feature vector used for state discovery, and let $S$ be the true latent state (equilibrium, thermal, etc.). Define the \emph{feature weakness} $\delta = I(\phi(X); S)$, the mutual information between the features and the true state; when $\delta$ is small, the features carry little information about which state a sample belongs to. We prove (Supplementary Information, Section~S2):

\begin{equation}
\mathrm{F1}_{\mathrm{SCX}} \;\leq\; \mathrm{F1}_{\mathrm{base}} \;+\; C_F\,\sqrt{\frac{\delta}{2}},
\label{eq:theorem2}
\end{equation}

where $\mathrm{F1}_{\mathrm{base}}$ is the F1 score of the best possible state-agnostic method (e.g., thresholding the loss) and $C_F$ is a constant depending on the minimum precision and recall of the detector (typically $C_F \leq 2$ when precision and recall exceed $0.1$, and $C_F \leq 1$ when both exceed $0.5$). The $\sqrt{\delta/2}$ factor follows from Pinsker's inequality: $TV(P, \tilde{P}) \leq \sqrt{\text{KL}(P\|\tilde{P})/2} = \sqrt{\delta/2}$. As $\delta \to 0$, the SCX F1 converges to that of the baseline---the multi-expert consistency test cannot improve over a simple loss threshold when the features are uninformative.

This is not a weakness of the method but an honest statement of its boundary conditions. Equation~\eqref{eq:theorem2} provides a practical diagnostic: before running SCX, estimate $\delta$ from the data; if $\delta$ is small relative to $\log K$ (where $K$ is the number of candidate states), the method will not meaningfully outperform simpler approaches. The DermaMNIST case in our cross-domain validation (Section~\ref{sec:cross_domain}) exemplifies this failure mode: 28$\times$28 thumbnail images processed through a shallow SimpleCNN produce feature representations with $\delta \approx 0.02$, and SCX indeed degrades to near-baseline performance (F1 = 0.101 vs baseline 0.105 at 10\% noise).

Crucially, the bound in Equation~\eqref{eq:theorem2} also tells us what to do about weak features: the remedy is not to tune SCX hyperparameters but to enrich the feature representation. Switching from SimpleCNN to a deeper ResNet-18 backbone on the same DermaMNIST data increases $\delta$ to $0.18$, and SCX F1 rises to $0.243$---still limited, but a meaningful improvement. For the AlN MLIP, replacing the 12-dimensional handcrafted descriptors with 182-dimensional ACE descriptors reduces the normalized weakness $\epsilon_\phi = \delta / \log K$ from $0.46$ to an estimated $0.12$, which is consistent with the observed F1 improvement from $0.45$ to $0.585$.

\subsection{Why This Problem Is Fundamentally Hard}
\label{sec:unidentifiability}

A skeptical reader might ask: why train multiple experts at all? Why not simply inspect each sample and determine whether its label is correct? The answer is subtle but important: without additional assumptions, the problem of distinguishing noise from hardness is provably unsolvable from observational data alone.

We prove this formally (Supplementary Information, Section~S3). Consider any data distribution $P_{X,Y,\tilde{Y}}$ where $X$ is the input, $Y$ is the true label, and $\tilde{Y}$ is the observed (possibly noisy) label. Any observable quantity---the loss, the expert predictions, the input features---is a function of $(X, \tilde{Y})$ alone. For any such function $f$, it is possible to construct two different data-generating processes:

\begin{enumerate}
    \item \textbf{Noise world}: $\tilde{Y}$ is a noisy corruption of $Y$, with noise rate $\rho(x)$ that depends on the input.
    \item \textbf{Hardness world}: $\tilde{Y} = Y$ (all labels are correct), but the Bayes-optimal error rate $\varepsilon(x)$ varies across the input space.
\end{enumerate}

These two worlds produce observationally identical distributions over $(X, \tilde{Y})$. No algorithm, no matter how sophisticated, can distinguish between them without access to information beyond what is contained in the observed data. In the language of statistical identifiability, the noise rate $\rho(x)$ and the hardness $\varepsilon(x)$ are not separately identifiable from the marginal distribution of $(X, \tilde{Y})$ alone.

This result---which we call the Unidentifiability Theorem---has a constructive corollary: it tells us exactly what additional structure is needed to break the ambiguity. Our six assumptions (A1--A6) are not arbitrary technical conditions; they are the minimal set required to make noise and hardness separately identifiable. Assumption A1 (expert training on disjoint data) ensures that expert errors are conditionally independent given the state, enabling the multiplicity-of-failures test. Assumption A4 (uniform independent noise) ensures that the noise rate does not vary with the input in a way that mimics hardness. Assumption A5 (state-homogeneous error rates) provides the calibration reference against which all-expert failures can be measured. Each assumption corresponds to a specific path out of the identifiability trap.

We are aware of no prior data cleaning method that acknowledges this theoretical boundary. Most approaches implicitly assume that noise is separable from difficulty through a loss threshold, a confidence score, or a disagreement measure---but none proves that such separation is possible, nor states the conditions under which it fails. The Unidentifiability Theorem reveals that this oversight is not a minor omission but a fundamental gap: without explicitly stated assumptions, noise detection methods are solving an ill-posed problem.

\subsection{Practical Implementation}
\label{sec:implementation}

The complete SCX pipeline proceeds as follows. Given a training dataset $\mathcal{D} = \{(x_i, y_i)\}_{i=1}^N$ and a choice of domain-appropriate features:

\begin{enumerate}
    \item \textbf{Feature extraction} (Layer 1): Compute domain-knowledge features $\phi(x_i)$ for each sample. For MLIP, these are ACE descriptors; for vision, CNN penultimate-layer activations from a pretrained backbone; for tabular data, the raw features.
    \item \textbf{Error-driven encoding} (Layer 2): Train $M_{\mathrm{probe}}$ probe models on random data subsets. For each sample, compute the probe error vector, then learn a projection matrix $W$ that maximizes the mutual information between the projected features and the probe error pattern.
    \item \textbf{State assignment}: Cluster the projected features into $K$ states using k-means with a cluster-balance constraint ensuring each state has sufficient statistical mass.
    \item \textbf{Expert training}: Train $M$ expert models on disjoint random subsets of $\mathcal{D}$. Experts share the model architecture but have independent initialization and training trajectories.
    \item \textbf{Noise detection}: For each sample $x$ in state $s$, compute the consistency score $C(x)$ as the fraction of experts that mispredict. Flag as noise if $C(x) > \tau(s)$, where $\tau(s) = \bar{e}_s + \kappa \sigma_s$ is set to the mean plus $\kappa$ standard deviations of the clean-sample error rate in state $s$, calibrated using Theorem~1.
    \item \textbf{Weak feature diagnosis}: Before interpreting results, estimate $\delta = \hat{I}(\phi(X); S)$ and compute the normalized weakness $\epsilon_\phi$. If $\epsilon_\phi > 0.5$, issue a warning that the detected noise labels may not be reliable, consistent with Theorem~2.
\end{enumerate}

The entire pipeline requires two $M$-model training passes: one for probe models (state discovery) and one for detection experts. In practice, the probe models can be low-epoch, low-capacity versions of the full model, keeping the total computational cost modest. For the AlN v3 dataset (534 frames, ACE descriptors), the complete pipeline completes in under 2 hours on a single CPU node; for CIFAR-10 (50,000 images, ResNet-18 backbone), it requires approximately 8 GPU-hours with $M = 20$ experts trained for 10 epochs each.

% End of theory_methods.tex
